Une station de radio utilise deux antennes verticales séparées de \qty{\v{D}}{m}, alimentées en phase à la fréquence de \qty{\v{f}}{MHz}. On étudie l'intensité reçue loin des antennes, dans le plan horizontal, en fonction de l'angle $\theta$ mesuré par rapport à la médiatrice du segment qui relie les antennes.

\begin{enumerate}[label=\alph*)]
    \item Déterminez la longueur d'onde des ondes émises.
    \item À quel angle se trouve le premier maximum d'intensité après celui de la médiatrice ?
    \item Combien de directions de maximum d'intensité y a-t-il entre $\theta = \qty{-90}{\degree}$ et $\theta = \qty{90}{\degree}$ ?
    \item À quel angle un auditeur ne capte-t-il presque rien pour la première fois ?
\end{enumerate}
